\documentclass[11pt,a4paper]{article}

% --- Packages ---
\usepackage[utf8]{inputenc}
\usepackage[margin=1in]{geometry}
\usepackage{amsmath, amsfonts, amssymb, amsthm}
\usepackage{xcolor}
\usepackage{mdframed} % For boxing the final results
\usepackage{titlesec} % For styled headings
\usepackage{hyperref}

% --- Custom Styling ---
\definecolor{navyblue}{RGB}{0, 0, 128}
\titleformat{\section}{\large\bfseries\color{navyblue}}{\thesection}{1em}{}[{\titlerule[0.8pt]}]
\titleformat{\subsection}{\normalsize\bfseries}{\thesubsection}{1em}{}

% --- Document Info ---
\title{\vspace{-2cm} \textbf{Black--Scholes Transformation Derivation}}
\author{
    \textbf{Parth Malav} \quad 
    \textbf{Shashank Kumar} \quad 
    \textbf{Kushagra Mani Tripathi} \\
    \small Project Contributors
}
\date{\today}

\begin{document}

\maketitle

\section{Introduction}
The Black--Scholes Partial Differential Equation (PDE) governs the price evolution of a European-style derivative. This report details the transformation of the Black--Scholes PDE into the standard heat equation.

The original PDE is given by:
\begin{equation}
\frac{\partial V}{\partial t}
+ \frac{1}{2}\sigma^2 S^2 \frac{\partial^2 V}{\partial S^2}
+ rS \frac{\partial V}{\partial S}
- rV = 0
\end{equation}

\section{Change of Variables}
To simplify the equation, we introduce a logarithmic change of spatial variables and a time-reversal transformation:

\begin{mdframed}[linecolor=gray!20, backgroundcolor=gray!5]
\begin{equation}
x = \ln\left(\frac{S}{K}\right), \quad \tau = \frac{1}{2}\sigma^2 (T - t)
\end{equation}
\end{mdframed}

We assume a solution of the form:
\begin{equation}
V(S,t) = e^{-\alpha x - \beta \tau} u(x,\tau)
\end{equation}

\section{Derivation of Partial Derivatives}

\subsection{Time Derivative}
Using the chain rule for $\tau(t)$:
\begin{equation}
\frac{\partial V}{\partial t} = \frac{\partial V}{\partial \tau} \frac{\partial \tau}{\partial t} = -\frac{\sigma^2}{2} e^{-\alpha x - \beta \tau} \left( -\beta u + \frac{\partial u}{\partial \tau} \right)
\end{equation}

\subsection{Spatial Derivatives}
First, the first-order derivative with respect to $S$:
\begin{equation}
\frac{\partial V}{\partial S} = \frac{\partial V}{\partial x} \frac{\partial x}{\partial S} = \frac{1}{S} e^{-\alpha x - \beta \tau} \left( -\alpha u + \frac{\partial u}{\partial x} \right)
\end{equation}

Next, the second-order derivative:
\begin{align}
\frac{\partial^2 V}{\partial S^2} &= \frac{\partial}{\partial S} \left( \frac{\partial V}{\partial S} \right) \notag \\
&= \frac{1}{S^2} e^{-\alpha x - \beta \tau} \left[ \alpha^2 u - (2\alpha+1)\frac{\partial u}{\partial x} + \frac{\partial^2 u}{\partial x^2} \right]
\end{align}

\section{Substitution and Coefficient Matching}
Substituting these expressions into the original PDE and dividing by common exponential factors, we group terms by $u$, $u_x$, and $u_\tau$:

\begin{equation}
-\frac{\partial u}{\partial \tau} + \frac{\partial^2 u}{\partial x^2} + \left[ \frac{2r}{\sigma^2} - (2\alpha + 1) \right] \frac{\partial u}{\partial x} + \left[ \alpha^2 + \left(\frac{2r}{\sigma^2} - 1\right)\alpha - \frac{2r}{\sigma^2} + \beta \right] u = 0
\end{equation}

\subsection{Determining $\alpha$ and $\beta$}
To reduce this to the \textbf{Heat Equation} ($u_\tau = u_{xx}$), we set the coefficients of $u_x$ and $u$ to zero:

\begin{itemize}
    \item \textbf{Solving for $\alpha$:} 
    $ \frac{2r}{\sigma^2} - (2\alpha + 1) = 0 \implies \alpha = \frac{r}{\sigma^2} - \frac{1}{2} $
    
    \item \textbf{Solving for $\beta$:} 
    Substituting $\alpha$ into the $u$ coefficient yields:
    $ \beta = \left( \frac{r}{\sigma^2} + \frac{1}{2} \right)^2 $
\end{itemize}

\section{Conclusion}
The Black--Scholes PDE has been successfully transformed into the dimensionless Heat Equation:

\begin{mdframed}[linecolor=navyblue, linewidth=1.5pt]
\begin{equation}
\frac{\partial u}{\partial \tau} = \frac{\partial^2 u}{\partial x^2}
\end{equation}
\end{mdframed}

This allows for the application of standard Green's function solutions to price options.

\end{document}